% Generated by publication/build_sections.py; do not edit by hand.
\paragraph{Informal verdict.}
\textbf{A --- Complete solution}, with \textbf{very high confidence}.

\paragraph{Lean coverage.}
The Lean development proves the printed normalized Bezout-pair theorem, its uniqueness, the universal binomial identity, the classification of all pairs, the worked example, and the constant-polynomial boundary case over an arbitrary field.

\subsection{Audited informal answer}
\subsubsection{1. Faithful restatement}

Let \(\mathbb K\) be a field. Let \(A,B\in\mathbb K[X]\) be coprime, and let \(U,V\in\mathbb K[X]\) satisfy

\[
\deg U<\deg B,
\qquad
AU+BV=1.
\]

For positive integers (m,n), determine explicitly the pair \((U_{m,n},V_{m,n})\in\mathbb K[X]^2\) such that

\[
\deg U_{m,n}<\deg B^n
\qquad\text{and}\qquad
A^mU_{m,n}+B^nV_{m,n}=1.
\]

This is exactly Q604 on the rendered PDF page. The notation \(d^\circ\) denotes polynomial degree. There is no apparent typo. The statement leaves implicit the standard convention that \(\mathbb K\) is a field; this is the setting used below. We use \(\deg 0=-\infty\). 

\medskip\hrule\medskip

\subsubsection{2. Explicit answer}

For integers \(r\ge 1\) and \(s\ge 1\), define

\[
S_{r,s}(T)
:=
\sum_{j=0}^{s-1}\binom{r+j-1}{j}T^j
\in \mathbb K[T],
\]

where the integer binomial coefficients are mapped into \(\mathbb K\) through the canonical homomorphism \(\mathbb Z\to\mathbb K\).

Set

\[
\widehat U_{m,n}
:=
U^mS_{m,n}(BV)
=
%
U^m\sum_{j=0}^{n-1}
\binom{m+j-1}{j}(BV)^j
\]

and

\[
\widehat V_{m,n}
:=
V^nS_{n,m}(AU)
=
%
V^n\sum_{i=0}^{m-1}
\binom{n+i-1}{i}(AU)^i.
\]

Perform the Euclidean division

\[
\widehat U_{m,n}=B^nQ_{m,n}+R_{m,n},
\qquad
\deg R_{m,n}<\deg B^n.
\]

Then the required pair is

\[
\boxed{
U_{m,n}=R_{m,n},
\qquad
V_{m,n}=\widehat V_{m,n}+A^mQ_{m,n}.
}
\]

Equivalently, writing \(\operatorname{rem}_{C}(P)\) for the Euclidean remainder of \(P\) modulo a nonzero polynomial \(C\),

\[
\boxed{
U_{m,n}
=
%
\operatorname{rem}_{B^n}
\left(
U^m
\sum_{j=0}^{n-1}
\binom{m+j-1}{j}(BV)^j
\right),
}
\]

and then

\[
\boxed{
V_{m,n}
=
%
\frac{1-A^mU_{m,n}}{B^n}.
}
\]

The quotient in the last formula is a polynomial, as proved below.

\medskip\hrule\medskip

\subsubsection{3. The universal binomial Bézout identity}

The formula rests on the following identity.

\paragraph{Lemma}

Let \(R\) be any commutative ring, and let \(x,y\in R\) satisfy \(x+y=1\). For all \(r,s\ge 1\),

\[
\boxed{
x^rS_{r,s}(y)+y^sS_{s,r}(x)=1.
}
\]

Thus explicitly,

\[
x^r
\sum_{j=0}^{s-1}\binom{r+j-1}{j}y^j
+
y^s
\sum_{i=0}^{r-1}\binom{s+i-1}{i}x^i
=1.
\]

\paragraph{Proof}

For convenience put \(S_{r,0}=0\). Pascal's identity gives, for \(r\ge 2\),

\[
S_{r,s}(T)-S_{r-1,s}(T)
=
%
T S_{r,s-1}(T).
\]

Indeed, for \(j\ge1\),

\[
\binom{r+j-1}{j}-\binom{r+j-2}{j}
=
%
\binom{r+j-2}{j-1}.
\]

Since \(x=1-y\),

\[
\begin{aligned}
S_{r-1,s}(y)-xS_{r,s}(y)
&=
S_{r-1,s}(y)-(1-y)S_{r,s}(y)\\
&=
S_{r-1,s}(y)-S_{r,s}(y)+yS_{r,s}(y)\\
&=
-yS_{r,s-1}(y)+yS_{r,s}(y)\\
&=
\binom{r+s-2}{s-1}y^s.
\end{aligned}
\]

Consequently,

\[
x^{r-1}S_{r-1,s}(y)-x^rS_{r,s}(y)
=
%
\binom{r+s-2}{s-1}x^{r-1}y^s.
\]

On the other hand,

\[
S_{s,r}(x)-S_{s,r-1}(x)
=
%
\binom{s+r-2}{r-1}x^{r-1}.
\]

Because

\[
\binom{r+s-2}{s-1}
=
%
\binom{r+s-2}{r-1},
\]

we obtain

\[
x^rS_{r,s}(y)+y^sS_{s,r}(x)
=
%
x^{r-1}S_{r-1,s}(y)+y^sS_{s,r-1}(x).
\]

Iterating this reduction in \(r\) gives

\[
x^rS_{r,s}(y)+y^sS_{s,r}(x)
=
%
xS_{1,s}(y)+y^s.
\]

Finally,

\[
S_{1,s}(y)=1+y+\cdots+y^{s-1},
\]

so

\[
xS_{1,s}(y)+y^s
=
%
= (1-y)(1+y+\cdots+y^{s-1})+y^s
%
1.
%
\]

This proves the lemma. Since it is an identity with integer coefficients, it remains valid in every characteristic. \(\square\)

\medskip\hrule\medskip

\subsubsection{4. Proof of the formula}

Apply the lemma with

\[
x=AU,\qquad y=BV.
\]

The given Bézout identity says precisely that \(x+y=1\). Therefore

\[
(AU)^mS_{m,n}(BV)
+
(BV)^nS_{n,m}(AU)
=
%
1.
%
\]

Factoring out \(A^m\) and \(B^n\) gives

\[
A^m
\left(
U^mS_{m,n}(BV)
\right)
+
B^n
\left(
V^nS_{n,m}(AU)
\right)
=
%
1,
\]

that is,

\[
A^m\widehat U_{m,n}+B^n\widehat V_{m,n}=1.
\tag{1}
\]

Thus \((\widehat U_{m,n},\widehat V_{m,n})\) is already an explicit Bézout pair for \(A^m\) and \(B^n\). It need not, however, satisfy the required degree condition.

Now divide

\[
\widehat U_{m,n}=B^nQ_{m,n}+R_{m,n},
\qquad
\deg R_{m,n}<\deg B^n.
\]

Substitution into \(1\) yields

\[
\begin{aligned}
1
&=
A^m(B^nQ_{m,n}+R_{m,n})
+
B^n\widehat V_{m,n}\\
&=
A^mR_{m,n}
+
B^n\bigl(\widehat V_{m,n}+A^mQ_{m,n}\bigr).
\end{aligned}
\]

Hence

\[
U_{m,n}=R_{m,n},
\qquad
V_{m,n}=\widehat V_{m,n}+A^mQ_{m,n}
\]

has all the required properties.

In particular,

\[
B^n\mid 1-A^mU_{m,n},
\]

so

\[
V_{m,n}=\frac{1-A^mU_{m,n}}{B^n}
\]

is indeed a polynomial.

\medskip\hrule\medskip

\subsubsection{5. Uniqueness and classification of all pairs}

Suppose that two pairs ((P,Q)) and ((P',Q')) satisfy

\[
A^mP+B^nQ=1,
\qquad
A^mP'+B^nQ'=1.
\]

Subtracting gives

\[
A^m(P-P')=-B^n(Q-Q').
\tag{2}
\]

Because \(1\) is itself a Bézout identity for \(A^m\) and \(B^n\), equation \(2\) implies that \(B^n\mid P-P'\). Indeed,

\[
\begin{aligned}
P-P'
&=
\bigl(A^m\widehat U_{m,n}+B^n\widehat V_{m,n}\bigr)(P-P')\\
&=
\widehat U_{m,n},A^m(P-P')
+
B^n\widehat V_{m,n}(P-P'),
\end{aligned}
\]

and both terms on the right are divisible by \(B^n\).

If \(B\) is nonconstant and

\[
\deg P,\deg P'<\deg B^n,
\]

then either \(P-P'=0\), or

\[
\deg(P-P')\ge \deg B^n,
\]

which is impossible. Thus \(P=P'\), and then \(2\) gives \(Q=Q'\).

If \(B\) is a nonzero constant, the degree condition is \(\deg P<0\), hence \(P=0\), so uniqueness is again immediate.

Therefore the normalized pair \((U_{m,n},V_{m,n})\) is unique.

Moreover, all Bézout pairs without the degree restriction are exactly

\[
\boxed{
\bigl(U_{m,n}+B^nT,;V_{m,n}-A^mT\bigr),
\qquad T\in\mathbb K[X].
}
\]

Thus the condition \(\deg U_{m,n}<\deg B^n\) selects precisely one pair from this affine family.

\medskip\hrule\medskip

\subsubsection{6. Symmetric remainder formula}

Assume for this subsection that \(A\) and \(B\) are both nonconstant. Then the normalized pair automatically satisfies

\[
\deg V_{m,n}<m\deg A.
\]

Indeed, if \(V_{m,n}\ne0\), the highest-degree terms in

\[
A^mU_{m,n}+B^nV_{m,n}=1
\]

must cancel, so

\[
m\deg A+\deg U_{m,n}
=
%
n\deg B+\deg V_{m,n}.
\]

Since \(\deg U_{m,n}<n\deg B\), it follows that

\[
\deg V_{m,n}<m\deg A.
\]

Now \(V_{m,n}\equiv\widehat V_{m,n}\pmod{A^m}\), because

\[
V_{m,n}=\widehat V_{m,n}+A^mQ_{m,n}.
\]

Consequently,

\[
\boxed{
V_{m,n}
=
%
\operatorname{rem}_{A^m}
\left(
V^n
\sum_{i=0}^{m-1}
\binom{n+i-1}{i}(AU)^i
\right).
}
\]

Thus, in the generic nonconstant case, the answer has the completely symmetric form

\[
\boxed{
\begin{aligned}
U_{m,n}
&=
\operatorname{rem}_{B^n}
\left(
U^mS_{m,n}(BV)
\right),\\[2mm]
V_{m,n}
&=
\operatorname{rem}_{A^m}
\left(
V^nS_{n,m}(AU)
\right).
\end{aligned}
}
\]

\medskip\hrule\medskip

\subsubsection{7. A directly degree-normalized \(B\)-adic formula}

The remainder modulo \(B^n\) can itself be expanded explicitly using only successive divisions by \(B\).

Assume \(\deg B=b\ge1\). Define

\[
\rho_B(P):=\operatorname{rem}_B(P),
\qquad
\delta_B(P):=\frac{P-\rho_B(P)}{B}.
\]

Thus

\[
P=\rho_B(P)+B\delta_B(P).
\]

Iterating gives the finite \(B\)-adic expansion

\[
P
=
%
\sum_{r\ge0}B^r\rho_B\bigl(\delta_B^r(P)\bigr).
\]

Since

\[
\widehat U_{m,n}
=
%
\sum_{k=0}^{n-1}
\binom{m+k-1}{k}
B^kU^mV^k,
\]

its terms of \(B\)-adic order strictly below \(n\) give

\[
\boxed{
U_{m,n}
=
%
\sum_{r=0}^{n-1}B^r C_r,
}
\]

where

\[
\boxed{
C_r
=
%
\sum_{k=0}^{r}
\binom{m+k-1}{k}
,\rho_B!\left(
\delta_B^{,r-k}(U^mV^k)
\right).
}
\]

Every \(C_r\) has degree \(<b\), so the displayed formula directly gives

\[
\deg U_{m,n}<nb=\deg B^n.
\]

This is equivalent to the single-remainder formula but may be preferable for an implementation based on \(B\)-adic digits.

\medskip\hrule\medskip

\subsubsection{8. Small and boundary cases}

\paragraph{Case \(m=n=1\)}

Here \(S_{1,1}=1\), so

\[
\widehat U_{1,1}=U,\qquad
\widehat V_{1,1}=V.
\]

Because \(\deg U<\deg B\), no reduction is needed:

\[
(U_{1,1},V_{1,1})=(U,V).
\]

\paragraph{Case \(n=1\)}

Since \(S_{m,1}=1\),

\[
\boxed{
U_{m,1}=\operatorname{rem}_B(U^m).
}
\]

The corresponding second coefficient is

\[
V_{m,1}=\frac{1-A^mU_{m,1}}B.
\]

Before normalization, the symmetric pair is

\[
\widehat U_{m,1}=U^m,
\qquad
\widehat V_{m,1}
=
%
V\sum_{i=0}^{m-1}(AU)^i.
\]

\paragraph{Case \(m=1\)}

Because \(S_{1,n}(T)=1+T+\cdots+T^{n-1}\),

\[
U_{1,n}
=
%
\operatorname{rem}_{B^n}
\left(
U\sum_{j=0}^{n-1}(BV)^j
\right),
\]

while

\[
\widehat V_{1,n}=V^n.
\]

The raw identity is simply

\[
A,U\sum_{j=0}^{n-1}(BV)^j+(BV)^n=1.
\]

\paragraph{Constant \(B\)}

If \(B\in\mathbb K^\times\), then \(\deg U<0\) forces \(U=0\), and \(AU+BV=1\) gives

\[
V=B^{-1}.
\]

The unique answer is therefore

\[
\boxed{
U_{m,n}=0,
\qquad
V_{m,n}=B^{-n}.
}
\]

This also covers the possible degenerate case \(A=0\).

\paragraph{Constant \(A\), nonconstant \(B\)}

If \(A\in\mathbb K^\times\) and \(B\) is nonconstant, the canonical initial pair is

\[
U=A^{-1},\qquad V=0.
\]

Hence

\[
\boxed{
U_{m,n}=A^{-m},
\qquad
V_{m,n}=0.
}
\]

\paragraph{Positive characteristic}

There are no exceptional characteristics. The coefficients

\[
\binom{m+j-1}{j}
\]

are simply reduced modulo \(\operatorname{char}\mathbb K\). The proof uses no factorial denominators and is an identity over \(\mathbb Z\).

\medskip\hrule\medskip

\subsubsection{9. Why the Euclidean reduction is genuinely necessary}

The raw pair \((\widehat U_{m,n},\widehat V_{m,n})\) need not satisfy the degree condition.

For example, over \(\mathbb Q\), take

\[
A=X^2,\qquad B=X+1,\qquad U=1,\qquad V=1-X.
\]

Then

\[
AU+BV=X^2+(X+1)(1-X)=1.
\]

For \(m=1,n=2\),

\[
\widehat U_{1,2}
=
%
= 1+BV
%
2-X^2.
\]

But

\[
\deg(2-X^2)=2=\deg (X+1)^2,
\]

so the strict degree condition fails. Dividing,

\[
2-X^2
=
%
-(X+1)^2+(2X+3).
\]

Thus

\[
U_{1,2}=2X+3.
\]

Also,

\[
\widehat V_{1,2}=V^2=(1-X)^2,
\qquad
Q_{1,2}=-1,
\]

so

\[
V_{1,2}
=
%
= (1-X)^2-X^2
%
1-2X.
\]

Indeed,

\[
X^2(2X+3)+(X+1)^2(1-2X)=1.
\]

\medskip\hrule\medskip

\subsubsection{10. Hypothesis audit}

The identity \(AU+BV=1\) already implies that \(A\) and \(B\) are coprime, so the explicit coprimality hypothesis is logically redundant once (U,V) are given.

The initial condition \(\deg U<\deg B\) is not needed for the construction: the same formula works starting from any Bézout pair \(AU+BV=1\). Its role is to make the given initial pair canonical.

The target condition

\[
\deg U_{m,n}<\deg B^n
\]

is essential for uniqueness. Without it, all solutions are obtained by adding

\[
(B^nT,-A^mT).
\]

Finally, the field hypothesis is used for unrestricted Euclidean division in \(\mathbb K[X]\). The universal binomial identity itself is valid over every commutative ring; the normalized construction also extends, for example, to coefficient rings over which division by the chosen \(B\) is available, such as when \(B\) is monic.
